% Section 3: the 2x2 original
\section{The $2\times2$ original: Duncan--Davis and JKST}\label{sec:2x2}

Source: PHASE1 \S0 (Appendix~\ref{app:phase1}), PHASE2a \S6.1, and the
original paper and package~\citep{jkst2019,duncan1953}.

\subsection{Duncan--Davis (the method of bounds)}

With two rows (Black, white; or group 1 and the rest) and a binary outcome, the
accounting identity in unit $i$ is $T_i=X_i\beta^b_i+(1-X_i)\beta^w_i$. Each
unit's rate for one group then lies in an interval:
\[
\beta^w_i\in\Big[\max\Big\{0,\tfrac{T_i-X_i}{1-X_i}\Big\},\ \min\Big\{1,\tfrac{T_i}{1-X_i}\Big\}\Big],
\]
and the same for $\beta^b_i$. Aggregating with the target weights gives the
Duncan--Davis (DD) bound on $B$ \citep{duncan1953}. In R$\times$C the per-cell
version is the Fréchet--Hoeffding bound
$n_{rc}\in[\max\{0,x_r+t_c-1\},\min\{x_r,t_c\}]$ (PHASE2a Prop.~8;
PHASE2b Prop.~D0).

\begin{plain}
DD uses nothing but arithmetic. It is always valid when A1 holds, never empty,
and often wide: in the BC $3\times3$ illustration of PHASE2a \S8 the DD bound
on loyalty is $[0.131,0.973]$.
\end{plain}

\subsection{JKST: linear contextual effects in $2\times2$}

JKST's Assumption~1 is that $(\beta^b_i,\beta^w_i,X_i,N_i)$ are iid with
$\E(\beta^w_i\mid X_i,N_i)=w_0+w_1X_i$ and
$\E(\beta^b_i\mid X_i,N_i)=b_0+b_1X_i$. The accounting identity then gives
\[
\E(T_i\mid X_i)=w_0+c_1X_i+d_1X_i^2 ,
\]
which identifies $(w_0,c_1,d_1)$ and leaves one free parameter, $w_1$.
\begin{itemize}[nosep]
\item \textbf{JKST Prop.~1} bounds $w_1$ through
$\E(L_i\mid X_i)\le w_0+w_1X_i\le \E(U_i\mid X_i)$ on a set $A$, where $L_i,U_i$
are the unit DD ends. It uses the conditional law of $T$ given $X$.
\item \textbf{JKST Prop.~2} relaxes Prop.~1 by Jensen's inequality and uses the
quadratic regression only. \emph{The package implements only Prop.~2}, with
$A=[\min X_i+10^{-5},\max X_i-10^{-5}]$, so the package bound is an outer set,
not the sharp one (PHASE1 \S0, fact~2).
\item \textbf{JKST Prop.~3} is the version with $A=[0,1]$ (extrapolation to pure
units).
\item \textbf{JKST Prop.~4 / App.~A (inference)}:
$\mathrm{CI}_x=[\hat B_L-xS_L,\ \hat B_U+xS_U]\cap DD$, with an HC1 sandwich
variance and the Popoviciu bound $\Var(\beta)\le1/4$ for the finite-population
term $S_1$. It needs tie-breaking (one active constraint) and covers the
parameter $B$ with probability at least $\Phi(x)$, in the spirit of
Imbens--Manski~\citep{imbensmanski2004}.
\item \textbf{App.~B}: the package reverts to DD when $\mathrm{CI}_0\cap DD$ is
empty; under the model $\mathrm{CI}_0\ne\emptyset$ in the large-$p$ limit.
\item \textbf{App.~C (covariates)}: covariates enter the intercept additively,
with a common slope; only $\tilde w_1$ is unidentified. Not in the package.
\end{itemize}

\begin{plain}
JKST's insight is that LCE leaves exactly one number unknown in $2\times2$: how
the curvature of the aggregate regression splits between the two groups'
contextual slopes. The requirement that both groups' implied rates be
probabilities at the two extreme compositions bounds that number, and hence $B$.
\end{plain}

\subsection{How the R$\times$C work contains JKST}

PHASE1 \S0 (fact~6) rewrote Prop.~2 in the homogeneous form of
Section~\ref{sec:setting} and checked numerically that it reproduces the
package's \texttt{generateBounds} to $10^{-16}$
(\path{docs/extension/check_2x2_equivalence.py}). PHASE2a Prop.~4
(Result~\ref{res:p2a-p4} below) proves the equivalence: the outer set $\ThO$ in
$2\times2$ is JKST Prop.~2 (with $A=[l,u]$) or Prop.~3 (with $A=\Delta^1$), the
dictionary is $K=c_1/2-w_1$, and $B^*(K)$ is JKST eq.~(16). PHASE2a Thm~1
(Result~\ref{res:p2a-t1}) confirms JKST's ``tightest bound'' claim for Prop.~1
when conditioning on $X$ only, and shows that conditioning on $(X,N)$ can be
strictly tighter.

\begin{revnote}
The Phases~1--3 referee asked that ``sharp'' never be applied to $\ThO$ without
``for the regression function'', and that the three sets, three hulls and three
targets be named in every statement. PHASE1 \S1.7 was added for this
(Appendix~\ref{app:revisions}, rows R-10, R-11 and R-33).
\end{revnote}
